\documentclass[11pt,letterpaper]{article}

\usepackage[margin=0.85in]{geometry}
\usepackage[T1]{fontenc}
\usepackage{lmodern}
\usepackage{microtype}
\usepackage{xcolor}
\usepackage{booktabs}
\usepackage{longtable}
\usepackage{array}
\usepackage{enumitem}
\usepackage{listings}
\usepackage{hyperref}
\usepackage{fancyhdr}

\definecolor{bladeblue}{HTML}{174A5B}
\definecolor{bladeorange}{HTML}{B85C2B}
\definecolor{codebg}{HTML}{F4F2EC}
\definecolor{codeframe}{HTML}{C8C3B8}

\hypersetup{
  colorlinks=true,
  linkcolor=bladeblue,
  urlcolor=bladeorange,
  citecolor=bladeblue,
  pdftitle={BLADE Complete Framework Manual},
  pdfauthor={BLADE Project}
}

\lstdefinestyle{blade}{
  basicstyle=\ttfamily\small,
  backgroundcolor=\color{codebg},
  frame=single,
  rulecolor=\color{codeframe},
  breaklines=true,
  columns=fullflexible,
  keepspaces=true,
  showstringspaces=false,
  tabsize=4
}
\lstset{style=blade}

\setlist{nosep}
\setlength{\parskip}{0.45em}
\setlength{\parindent}{0pt}
\setlength{\emergencystretch}{3em}
\pagestyle{fancy}
\fancyhf{}
\fancyhead[L]{BLADE Framework Manual}
\fancyhead[R]{\thepage}

\newcommand{\file}[1]{\nolinkurl{#1}}
\newcommand{\setting}[1]{\nolinkurl{#1}}
\newcommand{\stage}[1]{\textsf{#1}}

\title{\textbf{BLADE Complete Framework Manual}\\
\large Structure Generation, Thermodynamic Fitting, Phase Visualization,
Database Construction, and Oxidation Analysis}
\author{BLADE Project}
\date{July 18, 2026}

\begin{document}
\maketitle
\tableofcontents
\newpage

\section{Purpose and Scope}

BLADE is a general, high-throughput materials workflow that connects chemical
system enumeration, prototype-based structure generation, special
quasirandom structures (SQS), machine-learned interatomic-potential (MLIP)
relaxation, CALPHAD database fitting, thermodynamic visualization, reference
database construction, and environmental stability or oxidation analysis.

The example configurations use diboride-like and MAX-like structures to
demonstrate prototype definitions. A phase is defined by lattice vectors,
coordinates, variable and fixed sublattice labels, site species, and optional
fixed compositions. Fixed species such as B, C, or any other element are
specified through the user configuration.

The framework's practical novelty is the reproducible connection of:
\begin{enumerate}
  \item arbitrary prototype and sublattice definitions;
  \item high-throughput SQS generation through ATAT/mcsqs or SCRAPS;
  \item calculator-selectable structure relaxation and energy generation;
  \item automated fitting of per-system thermodynamic databases;
  \item binary through higher-order phase visualization;
  \item a shared energy database for environmental equilibrium calculations;
  \item static, animated, and exact-hover visual outputs; and
  \item cache-aware restart behavior across expensive stages.
\end{enumerate}

This is a workflow and software manual, not a validation claim for a specific
potential, fitted database, phase set, or scientific conclusion. Scientific
accuracy remains conditional on structures, energies, references, fitting
choices, solver assumptions, and convergence.

\section{Repository Layout}

\begin{longtable}{p{0.34\textwidth}p{0.60\textwidth}}
\toprule
Path & Purpose \\
\midrule
\endhead
\file{src/blade/tools/} & Composition, SQS, SCRAPS, cutoff, and TDB-generation classes. \\
\file{src/blade/analysis/} & Structural data extraction and thermodynamic/structure visualization. \\
\file{src/blade/oxidation/} & Database construction plus public oxidation calculation and visualization APIs. \\
\file{src/blade/oxidation/framework/} & General equilibrium, map, batch, single-composition, labeling, caching, and plotting implementation. \\
\file{examples/full_framework.py} & TOML-driven end-to-end orchestrator. \\
\file{examples/full_framework.toml} & Standard complete-workflow demonstration configuration. \\
\file{examples/full_framework_max.toml} & Multibasis demonstration configuration. \\
\file{examples/structures/} & Direct structure/SQS/TDB examples and complete prototype drivers. \\
\file{examples/oxidation/} & Direct database and oxidation examples. \\
\file{examples/extra/} & Optional refitting, comparison, and combined phase visualization utilities. \\
\file{Files/} & User-configurable persistent inputs, caches, databases, and outputs. \\
\file{uv.lock} & Reproducible uv dependency resolution. \\
\file{pixi.toml} and \file{pixi.lock} & Pixi environment and task definitions. \\
\bottomrule
\end{longtable}

\section{Installation}

\subsection{Python and Package Managers}

BLADE targets Python 3.12. Use either Pixi or uv.

\subsubsection{Pixi}
\begin{lstlisting}[language=bash]
cd /path/to/BLADE
pixi install
pixi run python -c "import blade; print(blade.__version__)"
pixi run check
\end{lstlisting}

\subsubsection{uv}
\begin{lstlisting}[language=bash]
cd /path/to/BLADE
uv sync --locked --extra workflow
uv run --locked --extra workflow python -c "import blade; print(blade.__version__)"
uv run --locked --extra workflow python examples/full_framework.py --check
\end{lstlisting}

The \setting{workflow} extra installs MaterialsFramework, pycalphad, mp-api,
and the configured MLIP dependency. The lock file should be committed and
retained with a run.

\subsection{External Programs}

\begin{longtable}{p{0.22\textwidth}p{0.70\textwidth}}
\toprule
Program & Requirement \\
\midrule
\endhead
ATAT & Supplies \file{sqs2tdb}, \file{corrdump}, and \file{mcsqs}; required by the mcsqs SQS/fitting path. \\
SCRAPS & Required only when \setting{tdb.sqs_method = "scraps"}. Configure \setting{paths.scraps_repo}. \\
FFmpeg & Required for MP4 movies and synchronized phase-grid video output. \\
Materials Project & Network access and an API key are required only for uncached reference downloads. \\
MLIP runtime & Calculator-dependent models, device support, and model files must be available. \\
\bottomrule
\end{longtable}

Store the Materials Project credential in the environment:
\begin{lstlisting}[language=bash]
export MP_API_KEY="your-key"
\end{lstlisting}
Do not commit credentials to TOML, source, logs, or documentation.

\section{End-to-End Orchestrator}

\subsection{Commands}

The default input is \file{examples/full_framework.toml}.
\begin{lstlisting}[language=bash]
# Validate files, settings, imports, and external dependencies.
pixi run python examples/full_framework.py --check

# Print enabled stages without executing calculations.
pixi run python examples/full_framework.py --dry-run

# Run the default configuration.
pixi run python examples/full_framework.py

# Run an alternate configuration.
pixi run python examples/full_framework.py examples/full_framework_max.toml
\end{lstlisting}

Equivalent uv commands prepend
\file{uv run --locked --extra workflow} instead of \file{pixi run}.

\subsection{Major Stage Switches}

The top of each TOML exposes the high-level controls:
\begin{lstlisting}
[stages]
sqs_generation = true
tdb_fitting = true
phase_diagrams = true
oxide_database = true
oxidation_graphs = true
\end{lstlisting}

The orchestrator executes these internal groups in order:
\begin{enumerate}
  \item \stage{tdb}: SQS generation and/or TDB fitting;
  \item \stage{phase\_visualization}: phase plots/GIFs and grid movies;
  \item \stage{database}: reference structures, energies, and phase database;
  \item \stage{oxidation}: batch or single-composition analysis.
\end{enumerate}

The two first booleans both route through the selected TDB driver. The runner
injects independent \file{run_sqs} and \file{run_tdb} values, so either half
can be run alone. The remaining booleans enable their corresponding stages.

\subsection{Configuration Injection}

\file{full_framework.py} parses the selected TOML with the Python standard
library. For legacy direct examples, it parses the target script's abstract
syntax tree, removes only top-level assignments whose names are supplied by
the orchestrator, and executes the script with injected values. Functions and
workflow code remain the same as direct execution. Database and oxidation
stages use class APIs directly.

\section{Complete TOML Reference}

\subsection{\texttt{[paths]}}

\begin{longtable}{p{0.32\textwidth}p{0.60\textwidth}}
\toprule
Setting & Meaning \\
\midrule
\endhead
\setting{blade_root} & Repository root used to find \file{src/} and \file{examples/}. \\
\setting{files_dir} & Persistent BLADE data/output root. \\
\setting{sqsdb_dir} & ATAT SQS database location used by the fitting tools. \\
\setting{scraps_repo} & SCRAPS checkout; used only by the SCRAPS backend. \\
\bottomrule
\end{longtable}

\subsection{\texttt{[tdb]}}

\begin{longtable}{p{0.35\textwidth}p{0.57\textwidth}}
\toprule
Setting & Meaning \\
\midrule
\endhead
\setting{sqs_method} & \file{"mcsqs"} or \file{"scraps"}. \\
\setting{prototype_driver} & \file{"standard"} or \file{"multibasis"}; selects the bundled general driver implementation. \\
\setting{level} & SQS composition level to generate/fit. \\
\setting{skip_existing_sqs} & Reuse completed SQS artifacts. \\
\setting{skip_existing_tdb} & Reuse existing TDB output. \\
\setting{refit_existing_tdb} & Refit from retained structures/energies even when SQS generation is skipped. \\
\setting{skip_existing_plots} & Reuse each plot independently when that output exists. \\
\setting{generate_gibbs_energy} & Enable per-binary Gibbs-energy plot. \\
\setting{generate_gibbs_mixing} & Enable per-binary Gibbs-mixing plot. \\
\setting{generate_phase_diagram} & Enable applicable phase-diagram plot. \\
\setting{generate_combined_phase_diagram} & Enable montage/combined diagram output. \\
\setting{generate_contcar_plots} & Enable combined structure images. \\
\setting{run_movie} & Enable relaxation movie generation. \\
\setting{mlip} & MaterialsFramework calculator registry name. \\
\setting{primary_elements}, \setting{secondary_elements} & Element pools used for system enumeration. \\
\setting{primary_min/max}, \setting{secondary_min/max} & Number of elements drawn from each pool per system. \\
\setting{liquid} & Include configured liquid treatment in fitting. \\
\setting{comps_folder} & Per-system output directory below \file{Files/}. \\
\bottomrule
\end{longtable}

\setting{[tdb_mlip_kwargs]} is passed to the selected calculator, for example
\setting{steps} and \setting{device}. Calculator-specific arguments are
allowed when supported by MaterialsFramework.

\subsection{\texttt{[tdb\_fit]}, \texttt{[tdb\_sqs]}, and \texttt{[tdb\_scraps]}}

\setting{[tdb_fit]} controls relaxation/fitting: \setting{fmax},
\setting{verbose}, \setting{t_min}, \setting{t_max}, \setting{sro},
\setting{bv}, \setting{phonon}, \setting{open_calphad}, and
\setting{track_trajectory}. Disabling trajectory tracking reduces storage.

\begin{sloppypar}
\setting{[tdb_sqs]} controls the mcsqs run time,
cutoff interpretation, pair/triplet/quadruplet cutoffs, correlation weights
\setting{wr}, \setting{wn}, and \setting{wd}, plus the parallel run count. Cutoff meaning
depends on \setting{cutoff_mode}; record the mode with numeric values.
\end{sloppypar}

\setting{[tdb_scraps]} controls SCRAPS ranks, automatic search budget,
maximum shell, and multibasis sublattice handling.

Each \setting{[[tdb_sqs_levels]]} has a numeric \setting{level}, one or more
\setting{compositions}, and the sublattice \setting{letter} list to which the
composition applies.

\subsection{\texttt{[phase]}}

\begin{longtable}{p{0.30\textwidth}p{0.62\textwidth}}
\toprule
Setting & Meaning \\
\midrule
\endhead
\setting{key} & Phase definition key referenced by the driver. \\
\setting{generator_name} & Base phase/prototype name. \\
\setting{a,b,c} and angles & Initial cell parameters and angles. \\
\setting{supercell_size} & Integer replication along the three lattice directions. \\
\setting{use_average_lattice} & Average active elements' lattice references for the phase. \\
\setting{vectors} & Three lattice vectors in the prototype convention. \\
\setting{coords} & Fractional coordinates followed by variable-site letters or fixed element symbols. \\
\bottomrule
\end{longtable}

When averaging is enabled, active elements come from the configured phase
site definitions or phase element pools. Prototype presets may be used when
available; otherwise periodic-table reference cells are used. Missing bulk
reference cells fall back to an isotropic covalent-radius estimate. Explicit
phase-level or configured element values can override inferred references.
Always record the resulting cell values because these are initial structure
estimates, not universal experimental constants.

\subsection{\texttt{[tdb\_inputs]}}

\setting{terms_in} and \setting{mult_in} are tables keyed by phase name,
optionally including system-size suffixes such as \file{PHASE_2} or
\file{PHASE_3}. Resolution selects the most specific applicable entry.
\setting{sqsgen_in} can similarly override generated SQS composition text.

Site tables define elements and behavior per sublattice:
\begin{lstlisting}
[tdb_inputs.sites.PHASE.a]
elements = ["A", "B", "C"]
constant = false

[tdb_inputs.sites.PHASE.b]
elements = ["D", "E"]
constant = true
composition = [0.4, 0.6]
\end{lstlisting}

A constant site preserves the listed element order and fixed composition
instead of permuting it. System-specific override tables can replace terms,
multiplicities, sites, SQS text, or fixed compositions for a selected chemical
system. Fractions on a fixed site must be numeric, match the element count,
and sum to one.

\subsection{\texttt{[phase\_plots]} and \texttt{[phase\_grid]}}

\setting{[phase_plots]} specifies input/output folders, plot temperature
range, missing-calculation policy, GIF settings, fixed species, variable-site
fraction, composition-grid density, convex-hull energy threshold, and
pressure. \setting{selected_systems} is inferred from configured TDB element
pools unless explicitly implemented by the calling context.

\setting{[phase_grid]} specifies the GIF directory, panel dimensions, optional
column count, output frame rate, background color, header/colorbar dimensions,
and temperature range. Its range and step must match GIF frame generation for
strict synchronization.

\subsection{\texttt{[database]}}

This section defines variable and fixed elements, oxidizing species, element
count limits, Materials Project stability filtering, duplicate formula
handling, POSCAR/energy cache behavior, MLIP name/label/folder, source
composition folders, four database substage switches, and
\setting{api_key_env}. Prefer naming an environment variable, normally
\file{MP_API_KEY}. An explicit \setting{api_key} value is also accepted.
For compatibility, an unset nondefault value in \setting{api_key_env} is
treated as a literal key, but committed credentials should be avoided.

\setting{[database_mlip_kwargs]} configures the reference-structure
calculator. \setting{[database_fallback_refs]} provides explicit elemental
energy references when the selected MLIP source does not provide one. Mixing
inconsistent reference methods changes formation energies and must be treated
as a scientific input choice.

\subsection{\texttt{[oxidation]} and Modes}

\setting{[oxidation]} selects \file{"batch"} or \file{"single"}, identifies an
optional fixed phase element and stoichiometry, chooses structure/database
subdirectories, configures region labels, and defines composition-slice axis
priority. Use no phase element and zero stoichiometry for a metal-only phase.

\setting{[oxidation_batch]} independently enables oxidation scans,
chemical-potential/composition maps, onset calculations, 3-D onset surfaces,
animations, composition slices, and per-composition slice muO--T maps.
The \setting{run_calculations} and \setting{run_plots} switches control two
global execution passes. When both are true, every selected system is
calculated and cached before plotting begins. Calculation-only mode disables
all figures and animations; plot-only mode performs no LP solves and reports
any missing required cache. A batch with failed systems prints the complete
failure summary and exits before a plotting pass can begin.
\setting{elements} is an optional allowed metal pool; an empty list includes
all systems, while a populated list skips systems containing any other metal.
The section also defines cache switches; temperature, chemical-potential, and
composition ranges; slice ratios; onset composition resolution; and onset
threshold. \setting{slice_muT_comp_step} controls only the fixed compositions
used for per-composition muO--T maps; it does not change muO--x slices or onset
composition grids. Batch muO--T maps are written only within composition slices.
Their region keys contain phase-fraction ranges, and their HTML files report
exact pointwise phase fractions on hover.
For each composition and temperature, an onset chemical-potential sweep stops
at the first infeasible assemblage. The onset table records both the last
feasible and first infeasible chemical potential. If no onset occurs over an
otherwise valid sweep, downstream AUC ranking uses the last feasible value,
not the configured upper chemical-potential bound.
The same sweep integrates raw LP fractions for the parent multicomponent phase
and for oxygen-containing phases. Parent-phase composition may vary or split
across multiple equilibrium pseudophases; all such fractions count while each
configured variable element remains above 0.01. End members missing an element
do not count toward the parent-phase integral. The private ranking utility then
integrates both quantities over temperature and plots parent-phase retention
against oxide formation.

\setting{[oxidation_single]} defines one system, metal ordering, normalized
initial composition, temperature and chemical-potential grids, scan
temperatures, composition coordinate range, cache setting, component-display
setting, Redlich--Kister order, and secondary map step. Its
\setting{run_calculations} and \setting{run_plots} switches use the same
calculation-first behavior as batch mode.

Ranges are three-value inclusive specifications \file{[start, stop, step]}.
The runner validates positive steps and materializes them with a small
floating-point endpoint allowance.

\section{Structure and TDB Calculations}

\subsection{Composition Enumeration}

\file{BladeCompositions} combines primary and secondary pools subject to
minimum and maximum counts. System identity is based on valid, unique chemical
symbols. The order-independent system key is generated from sorted element
symbols, while explicit sublattice order is retained where thermodynamically
meaningful.

\subsection{SQS Generation}

\file{BladeSQS} creates prototype and SQS composition inputs, asks ATAT to
materialize level directories, computes correlations, and launches one or
more mcsqs searches. Candidate quality is represented by retained correlation
outputs such as \file{bestcorr.out}. Pure-species directories may not need an
mcsqs search; their treatment is separate from a missing failed search.

\file{ScrapsSQSGen} provides the alternate SCRAPS path. Backend selection
should not change the downstream API, but it can change the selected
structures and therefore all fitted outputs.

\subsection{MLIP Relaxation and TDB Fit}

\file{BladeTDBGen} is configured without side effects. Calling \file{fit()}
iterates systems and phase data, prepares fitting directories, invokes the
MaterialsFramework fitting path, and writes TDBs. Calculator choice and
calculator arguments are forwarded explicitly; no calculator name is
hard-coded in the generalized class.

Refitting is safe with respect to avoiding a new MLIP calculation only when
the retained structures and energy files are valid and the intended changes
are limited to fitting/model inputs. It can change TDB coefficients and
thermodynamic predictions, as expected, while leaving source energies fixed.

\section{Phase Visualization}

\file{examples/extra/phase_visualization.py} is the single implementation for
phase plotting and synchronized grid compilation. Direct execution uses:
\begin{lstlisting}[language=Python]
run_phase_plots = True
run_phase_grid = True
\end{lstlisting}

The switches allow plot-only, grid-only, both, or neither. Plotting discovers
top-level TDBs under selected system folders. Binary systems receive Gibbs
energy, Gibbs mixing, and applicable phase-diagram plots. Higher-order systems
use a simplex composition grid, compute mixing free energy, identify the lower
convex hull, classify points above the hull threshold, and project the
N-component simplex to a regular polygon for display. Temperature frames can
be cached and assembled into GIFs.

Grid compilation discovers clean phase-evolution GIFs, groups them by parsed
element count, resizes panels, synchronizes frames to the configured
temperature array, adds temperature/colorbar context, and streams raw frames
to FFmpeg. If no GIF exists, grid-only mode reports the condition and returns
normally.

\section{Oxidation Database Workflow}

\subsection{Composition List}

\file{OxideCompositions} builds chemical systems from variable pools, fixed
phase species, and the oxidizing species. Inclusion flags control base,
oxidized, fixed-species, and no-oxygen combinations. It writes a composition
workbook used by reference acquisition.

\subsection{Reference Structures and Energies}

The Materials Project stage searches each chemical system, optionally limits
results to stable entries, skips formulas already represented by BLADE data,
and stores POSCARs with metadata. The MLIP stage uses one configured
calculator to relax/reference these structures and writes total and per-atom
energies.

\subsection{Database Assembly}

\file{OxideDatabase} scans BLADE structures and TDB outputs, combines MLIP
sources and fallback elemental references, writes unified and per-system
energy tables, and collects the files needed by equilibrium analysis into
\file{Files/system_structures/<system>/}. System detection derives from
configured directories and formulas; no phase element should be appended
implicitly to a system name.

\section{Oxidation Equilibrium and Maps}

\subsection{Public API}

\file{OxidationConfig} carries shared paths, phase species/stoichiometry,
labels, and slice priorities. \file{OxidationCalculator} method
\file{run_batch()} discovers
eligible systems and launches generalized analysis. \file{single()} returns a
fixed-composition analyzer for the requested system, metals, and composition.
\file{OxidationVisualizer} exposes visualization behavior separately from
calculation setup.

\subsection{Equilibrium Model}

At each temperature, oxygen chemical potential, and initial composition, the
framework constructs the feasible phase set and solves the phase-fraction
optimization subject to element mass balance and nonnegative fractions.
Consecutive chemical-potential solves are neighboring points in a scan; their
problem structure and cached thermodynamic data can be reused even though each
point represents a distinct equilibrium condition.

The configured initial composition belongs to the point being calculated. In
composition maps, each grid point is a different initial composition. A guide
line should appear on an onset plot only when the complete diagram is defined
relative to that initial composition.

\subsection{Region Identity and Labels}

Region separation is based on the set of displayed compounds after the
configured composition-component cutoff is applied. Fraction-range changes
alone do not create a separate region. If a component enters or leaves the
displayed phase composition at the cutoff, the displayed compound identity
changes and a new region may form.

Static keys summarize phase-fraction and phase-composition ranges. Endpoint
values near zero or one can be represented as 0 or 1 according to configured
labeling rules; small composition components can be omitted or included
according to the relevant option. Exact HTML hover text is generated from the
underlying point, not the summarized region key. Values below 0.01 retain
additional significant decimal places in hover text; ordinary values are
shown compactly.

\subsection{Outputs}

Depending on switches and system size, outputs include:
\begin{itemize}
  \item oxidation scan curves and per-temperature folders;
  \item oxygen-chemical-potential/composition assemblage maps;
  \item per-composition oxygen-chemical-potential/temperature slice maps;
  \item miscibility-gap panels and composition-slice maps;
  \item onset curves, AUC summaries, and 3-D onset surfaces;
  \item equilibrium phase-composition 3-D plots;
  \item GIF/MP4 animations;
  \item interactive HTML plots with exact point hover values;
  \item per-composition tables and compact NumPy cache files.
\end{itemize}

Composition-slice output uses a selected axis and configured remainder ratios.
Axis priority selects the first applicable element; it avoids generating every
permutation when only one ordered axis family is desired. Static muO--T maps
are stored in each slice's \file{muO_T_plots/} directory; GIF/MP4 and other
aggregate slice products remain at the slice root.

\section{Caching, Restarting, and Storage}

\begin{longtable}{p{0.30\textwidth}p{0.62\textwidth}}
\toprule
Artifact/setting & Reuse behavior \\
\midrule
\endhead
\file{bestcorr.out} and SQS outputs & Allow completed SQS searches to be skipped. Validate parseability; existence alone may not prove a valid result. \\
\setting{skip_existing_tdb} & Avoids a complete existing TDB fit unless refit is requested. \\
\setting{refit_existing_tdb} & Reuses relaxed energies but refreshes fitting inputs and TDB parameters. \\
\setting{skip_existing_plots} & Checks individual enabled plot files so missing products can be generated independently. \\
\setting{skip_if_tables_exist} & Reuses matching equilibrium tables instead of repeating solves. \\
\setting{skip_if_analysis_exists} & Reuses derived oxidation analysis products. \\
\file{.npz} caches & Compact numerical data used for redraw/interactive output. \\
Trajectory setting & Disable trajectory generation to reduce storage if no relaxation movie is needed. \\
\bottomrule
\end{longtable}

Do not reuse a cache after changing solver grids, source TDBs, phase energies,
element ordering, fixed species, initial compositions, model terms, or other
inputs that affect the cached quantity. Removing only a PNG or HTML file
should not require recalculation when the corresponding numerical cache is
intact.

\section{Every Example Script}

\begin{longtable}{p{0.35\textwidth}p{0.57\textwidth}}
\toprule
Script & Role and direct-run behavior \\
\midrule
\endhead
\file{examples/full_framework.py} & TOML orchestrator; defaults to the adjacent standard TOML. Supports \file{--check} and \file{--dry-run}. \\
\file{structures/01_compositions.py} & Standalone composition enumeration; no generated input required. \\
\file{structures/02_sqs_generation.py} & Standalone ATAT SQS generation from in-file prototype settings. \\
\file{structures/03_tdb_generation.py} & Standalone TDB fitting from existing SQS data. \\
\file{structures/04_visualization.py} & Standalone examples for structures, energies, phase diagrams, and movies; skips some absent inputs. \\
\file{structures/05_data_analysis.py} & Standalone scan of structure/energy directories into tabular data. \\
\file{structures/tdb_gen_hedb.py} & Complete standard/mcsqs prototype demonstration and orchestrator driver. \\
\file{structures/tdb_gen_hedb_scraps.py} & Complete standard/SCRAPS demonstration and driver. \\
\file{structures/tdb_gen_max.py} & Complete multibasis/mcsqs demonstration and driver. \\
\file{structures/tdb_gen_max_scraps.py} & Complete multibasis/SCRAPS demonstration and driver. \\
\file{oxidation/01_compositions.py} & Standalone oxidation-system enumeration. \\
\file{oxidation/02_poscars.py} & Standalone Materials Project acquisition; requires the step-01 workbook and a key. \\
\file{oxidation/03_energy.py} & Standalone reference-structure MLIP relaxation. \\
\file{oxidation/04_database.py} & Standalone database assembly and structure collection. \\
\file{oxidation/database_framework.py} & Resumable class-based combination of oxidation steps 01--04 using in-file settings. \\
\file{oxidation/oxidation_framework.py} & Standalone batch/single oxidation analysis using in-file settings. \\
\file{extra/phase_visualization.py} & Combined phase plotting/GIF and grid-video implementation with two booleans. \\
\file{extra/compare_energy.py} & Energy-source parity, residual, and tabular comparison. \\
\file{extra/refit_tdb.py} & Refit TDB model terms from existing structures and energies. \\
\file{extra/tdb_gen_alloy.py} & Complete single-sublattice alloy demonstration. \\
\file{extra/tdb_gen_hedb_quaternary.py} & Complete higher-order prototype demonstration. \\
\bottomrule
\end{longtable}

Every example Python file is directly executable, but later workflow steps
require the artifacts or services stated above. Library files under
\file{src/blade/} are imported APIs and are not intended as standalone
commands.

\section{Library Class and Module Reference}

\subsection{\texttt{blade.tools}}

\begin{itemize}
  \item \file{blade_compositions.py}: \file{BladeCompositions} generates
  constrained chemical systems.
  \item \file{blade_sqsgen.py}: \file{BladeSQS} creates ATAT SQS inputs and
  executes SQS generation.
  \item \file{blade_scraps_gen.py}: \file{ScrapsSQSGen} provides the SCRAPS
  backend.
  \item \file{blade_tdb_gen.py}: \file{BladeTDBGen} prepares and runs
  structure-energy/TDB fitting.
  \item \file{blade_cutoff.py}: neighbor/correlation cutoff support.
\end{itemize}

\subsection{\texttt{blade.analysis}}

\begin{itemize}
  \item \file{blade_data.py}: \file{BLADEData} extracts structural and energy
  information into data tables.
  \item \file{blade_visual.py}: \file{BladeVisualizer} produces structure
  images, Gibbs curves, phase diagrams, composites, and relaxation movies.
\end{itemize}

\subsection{\texttt{blade.oxidation}}

\begin{itemize}
  \item \file{compositions.py}: \file{OxideCompositions}.
  \item \file{database.py}: \file{OxideDatabase}.
  \item \file{oxidation_calculations.py}: public \file{OxidationConfig} and
  \file{OxidationCalculator}.
  \item \file{oxidation_visualization.py}: public
  \file{OxidationVisualizer}.
\end{itemize}

\subsection{\texttt{blade.oxidation.framework}}

\begin{longtable}{p{0.30\textwidth}p{0.62\textwidth}}
\toprule
Module & Responsibility \\
\midrule
\endhead
\file{config.py} & Shared analysis options and path/range configuration. \\
\file{system_config.py} & System discovery, metal detection, phase/TDB path assembly. \\
\file{thermodynamics.py} & Phase thermodynamic data and equilibrium/optimization support. \\
\file{compiler.py} & Input-table and phase-data compilation. \\
\file{analyzer.py} & Main analysis orchestration, region construction, labels, and plot dispatch. \\
\file{maps.py} & Composition and chemical-potential map calculations. \\
\file{single.py} & Fixed-composition workflow. \\
\file{batch.py} and \file{_batch_internal.py} & Multi-system batch workflow. \\
\file{_run.py} & Shared run entry behavior. \\
\file{utils.py} & Parsing, formatting, and numerical utilities. \\
\bottomrule
\end{longtable}

\section{Reproducible Execution Records}

For every scientific run, archive:
\begin{enumerate}
  \item the exact TOML or direct script;
  \item BLADE commit hash and uncommitted diff;
  \item \file{uv.lock} or \file{pixi.lock};
  \item Python, external binary, MLIP, and model versions;
  \item element pools/order, phase sites, cell, supercell, and fixed species;
  \item SQS backend, levels, cutoffs, budgets, seeds where available, and
  selected correlation output;
  \item relaxed structures, energies, reference energies, and fitting files;
  \item generated TDBs and phase database inventory;
  \item temperature, chemical-potential, and composition grids;
  \item cache settings and a manifest of reused artifacts;
  \item output tables, NumPy caches, and figures.
\end{enumerate}

Recommended command record:
\begin{lstlisting}[language=bash]
git rev-parse HEAD
git status --short
pixi run python examples/full_framework.py --check
pixi run python examples/full_framework.py --dry-run
pixi run python examples/full_framework.py 2>&1 | tee Files/run.log
\end{lstlisting}

\section{Quality Checks}

\begin{lstlisting}[language=bash]
# Repository lint and format verification
pixi run check

# Explicit Python compilation
pixi run python -m compileall -q src examples

# Configuration/dependency validation
pixi run python examples/full_framework.py --check

# Stage-routing validation without calculations
pixi run python examples/full_framework.py --dry-run

# uv lock and uv environment validation
uv lock --check
uv run --locked --extra workflow python examples/full_framework.py --check
\end{lstlisting}

To format edited Python:
\begin{lstlisting}[language=bash]
pixi run isort src examples
pixi run black src examples
pixi run ruff check src examples
\end{lstlisting}

\section{Troubleshooting}

\begin{longtable}{p{0.33\textwidth}p{0.59\textwidth}}
\toprule
Symptom & Interpretation and action \\
\midrule
\endhead
\file{No such file or directory: .../grace} & A calculator-specific folder/path was selected by configuration or old generated data. Set the calculator and all base paths in the direct script/TOML; generalized library code should not infer a calculator-named folder. \\
\file{metal detection ... positional argument} & Caller/method signature mismatch. Validate with the current package and runner versions; do not mix copied framework revisions. \\
\file{expected 2+ metals} & System detection found fewer variable species after fixed species were removed. Verify phase coordinates, fixed-element configuration, and directory/formula names. \\
\file{No bestcorr.out} & Pure species may legitimately skip mcsqs; otherwise the SQS search is incomplete. \\
\file{Could not parse objective} & The file exists but is not a valid completed correlation result. Do not treat it as a reusable completed SQS. \\
\file{No TDB files found} & TDB fitting did not run, selected output folder differs, or system filtering excludes the files. Check stage switches and \setting{comps_folder}. \\
\file{expected two metals} & A binary-only Gibbs/phase plot received a higher-order system. Higher-order phase-map logic should handle those systems instead. \\
\file{need at least one array}\newline
\file{to concatenate} & pycalphad found no tie-line data for the requested binary plot. Treat as a plot-specific failure; retained energy plots/TDB may still be valid. \\
\file{No clean GIFs found} & Grid compilation has no cached animation input. Enable phase plotting/GIF creation or point \setting{phase_diagrams_folder} to existing GIFs. \\
\file{TOMLDecodeError: Cannot overwrite a value} & The same key was assigned twice in one TOML scope. Remove or rename the duplicate assignment. \\
Plots recalculate unexpectedly & Preserve equilibrium tables/\file{.npz}, enable cache switches, and verify that grid/settings and output paths match the original run. \\
Missing plot families & Check both major stage switches and plot-specific booleans; a cache skip should be evaluated per output, not per entire plot group. \\
\bottomrule
\end{longtable}

\section{Extending to a New Material System}

\begin{enumerate}
  \item Copy a TOML to a new, versioned configuration.
  \item Define primary/secondary element pools and desired system sizes.
  \item Define the phase cell, vectors, coordinates, and supercell.
  \item Mark variable sites with site letters and fixed sites with actual
  species; do not add a demonstration fixed element unless it is physical.
  \item Define site element lists, constants, and compositions.
  \item Define SQS levels and backend settings.
  \item Set calculator and relaxation/fitting settings.
  \item Run \file{--check}, then SQS only.
  \item Inspect structures and correlation objectives.
  \item Run TDB fitting/refitting and inspect fit/phase outputs.
  \item Define database species/reference sources and construct the collected
  system database.
  \item Define oxidation phase species, grids, mode, and output switches.
  \item Run a reduced validation grid before the full unchanged-resolution
  production run.
  \item Archive all provenance listed above.
\end{enumerate}

\section{Standalone Versus Configured Execution}

\file{phase_visualization.py}, numbered examples, complete prototype drivers,
\file{database_framework.py}, and \file{oxidation_framework.py} can run with
their in-file settings and no TOML. \file{full_framework.py} is intentionally
TOML-driven, but its argument is optional because it defaults to
\file{examples/full_framework.toml}. Package modules are intended to be
imported into user code:
\begin{lstlisting}[language=Python]
from blade.tools import BladeCompositions, BladeSQS, BladeTDBGen
from blade.analysis import BLADEData, BladeVisualizer
from blade.oxidation import OxidationCalculator, OxidationConfig
\end{lstlisting}

The README files under \file{examples/} provide the shortest retraceable run
sequence. This manual provides the full reference and should be updated when
entry points, settings, output layout, or scientific assumptions change.

\end{document}
